\documentclass[12pt]{article}

% -----------------------------------------------------------------------
%  Packages
% -----------------------------------------------------------------------
\usepackage[margin=1in]{geometry}
\usepackage{amsmath,amssymb,amsthm}
\usepackage{mathtools}
\usepackage{bm}
\usepackage{booktabs}
\usepackage{hyperref}
\usepackage[numbers,sort&compress]{natbib}
\usepackage{microtype}
\usepackage{xcolor}
\usepackage{enumitem}
\usepackage{csquotes}
\usepackage[ruled,vlined,linesnumbered]{algorithm2e}

% -----------------------------------------------------------------------
%  Math macros (mirrors main.tex)
% -----------------------------------------------------------------------
\newcommand{\cS}{\mathcal{S}}
\newcommand{\cA}{\mathcal{A}}
\newcommand{\cO}{\mathcal{O}}
\newcommand{\cP}{\mathcal{P}}
\newcommand{\cZ}{\mathcal{Z}}
\newcommand{\cL}{\mathcal{L}}
\newcommand{\cI}{\mathcal{I}}
\newcommand{\cM}{\mathcal{M}}
\newcommand{\E}{\mathbb{E}}
\newcommand{\R}{\mathbb{R}}
\newcommand{\given}{\,\vert\,}
\newcommand{\argmax}{\operatorname{arg\,max}}

% algorithm2e style
\SetKwInput{KwInput}{Input}
\SetKwInput{KwOutput}{Output}
\SetKwInput{KwParam}{Parameters}
\SetKw{KwAnd}{and}
\SetKw{KwOr}{or}
\SetKwComment{Comment}{$\triangleright$\ }{}

% -----------------------------------------------------------------------
%  Metadata
% -----------------------------------------------------------------------
\title{\textbf{Algorithmic Specification for the\\
Intent-Constrained MDP Framework}\\[6pt]
\large Companion to: \emph{Mathematical Modeling of Autonomous Decision Agents
Using Markov Decision Processes with Intent Constraints}}

\author{Parag Kulkarni\\[4pt]
  \small Opus Technologies \quad|\quad Lovely Professional University\\
  \small \texttt{parag.kulkarni@opustechglobal.com}
}

\date{May 2026}

% -----------------------------------------------------------------------
\begin{document}
\maketitle

\begin{abstract}
This document provides formal algorithmic specifications for implementing
the Intent-Constrained Markov Decision Process (ICMDP) framework in the
context of Agentic Commerce. Seven algorithms are presented: LLM-based
intent parsing, intent-constrained value iteration and policy iteration,
masked Q-learning, masked policy gradient, real-time feasibility monitoring,
and the full end-to-end Agentic Commerce pipeline integrating all components.
Each algorithm is grounded in the formal definitions and equations of the
companion paper.
\end{abstract}

\tableofcontents
\newpage

% -----------------------------------------------------------------------
\section{Overview}
% -----------------------------------------------------------------------

The ICMDP framework requires three interoperating components at runtime:
\begin{enumerate}[leftmargin=*, label=(\roman*)]
  \item An \textbf{intent parser} that maps natural language instructions
  $l \in \cL$ and environment state $s \in \cS$ to an admissible action set
  $\cA_\cI(s) \subseteq \cA$.
  \item A \textbf{decision solver} that finds the optimal policy
  $\pi_\cI^*$ over the intent-restricted action space.
  \item A \textbf{feasibility monitor} that detects and handles the case
  $\cA_\cI(s) = \emptyset$ before the solver is invoked.
\end{enumerate}

The algorithms in Sections~\ref{sec:intent-parser}--\ref{sec:feasibility}
specify each component in isolation. Section~\ref{sec:pipeline} composes
them into the full Agentic Commerce execution pipeline.

% -----------------------------------------------------------------------
\section{Intent Parsing}
\label{sec:intent-parser}
% -----------------------------------------------------------------------

The intent parser $f_\theta : \cL \times \cS \to 2^{\cA}$ translates a
natural language instruction into a formal admissibility function. In
practice it is implemented as a structured LLM prompt that returns a set of
predicate constraints which are then evaluated against each candidate action.

\begin{algorithm}[H]
\caption{LLM Intent Parser}
\label{alg:intent-parser}
\KwInput{Natural language instruction $l \in \cL$;\;
         current state $s \in \cS$;\;
         full action set $\cA$;\;
         LLM with parameters $\theta$}
\KwOutput{Intent-admissible action set $\cA_\cI(s)$}

\BlankLine
$\text{prompt} \leftarrow \textsc{BuildPrompt}(l,\, s,\, \cA)$
\Comment*[r]{structured prompt requesting predicate constraints}

$\mathbf{c} \leftarrow \textsc{LLMCall}_\theta(\text{prompt})$
\Comment*[r]{$\mathbf{c} = \{c_1, \ldots, c_k\}$ formal predicates over actions}

$\cA_\cI(s) \leftarrow \emptyset$\;

\ForEach{$a \in \cA$}{
  \If{$c_j(s, a) = \mathbf{true}$ \KwAnd for all $j = 1,\ldots,k$}{
    $\cA_\cI(s) \leftarrow \cA_\cI(s) \cup \{a\}$\;
  }
}

\BlankLine
\textbf{return} $\cA_\cI(s)$\;
\end{algorithm}

\paragraph{Agentic Commerce instantiation.}
For the procurement scenario, the instruction
\enquote{\emph{Purchase 100 units from a carbon-neutral supplier, spending no
more than \$500, within 60 seconds}}, the parser produces three predicates:
$c_1$: $\text{eco}(v(a)) = 1$;\;
$c_2$: $p(a) \leq 500$;\;
$c_3$: $t_{\text{elapsed}} \leq 60\,\text{s}$.
Any action violating at least one predicate is excluded from $\cA_\cI(s)$.

% -----------------------------------------------------------------------
\section{Feasibility Monitoring}
\label{sec:feasibility}
% -----------------------------------------------------------------------

Before invoking any solver, the system must verify that the admissible set
is non-empty in every reachable state. Algorithm~\ref{alg:feasibility}
performs this check and escalates to constraint relaxation or user
re-specification when necessary.

\begin{algorithm}[H]
\caption{Feasibility Monitor}
\label{alg:feasibility}
\KwInput{Current state $s$;\;
         admissible set $\cA_\cI(s)$;\;
         relaxation budget $\delta_{\max}$}
\KwOutput{Feasible admissible set $\cA_\cI(s)$ or \textsc{Infeasible} signal}

\BlankLine
\If{$\cA_\cI(s) \neq \emptyset$}{
  \textbf{return} $\cA_\cI(s)$\;
}

\Comment{Attempt progressive constraint relaxation}
$\delta \leftarrow 0$\;
\While{$\delta \leq \delta_{\max}$}{
  $\cA_\cI^\delta(s) \leftarrow \textsc{RelaxConstraints}(\cI, s, \delta)$\;
  \If{$\cA_\cI^\delta(s) \neq \emptyset$}{
    $\textsc{NotifyUser}(\text{"Intent constraints relaxed by } \delta \text{"}$)\;
    \textbf{return} $\cA_\cI^\delta(s)$\;
  }
  $\delta \leftarrow \delta + \Delta\delta$\;
}

\Comment{Relaxation budget exhausted — surface to user}
$\textsc{ConstraintViolationWarning}(s, \cI)$\;
\textbf{return} \textsc{Infeasible}\;
\end{algorithm}

% -----------------------------------------------------------------------
\section{Planning Algorithms}
\label{sec:planning}
% -----------------------------------------------------------------------

\subsection{Intent-Constrained Value Iteration}

Algorithm~\ref{alg:icvi} is a direct adaptation of classical value
iteration \citep{puterman1994} in which the Bellman backup operator is
restricted to $\cA_\cI(s)$ at every state.

\begin{algorithm}[H]
\caption{Intent-Constrained Value Iteration (ICVI)}
\label{alg:icvi}
\KwInput{ICMDP $\cM_I = (\cS, \cA, \cP, R, \gamma, \cI)$;\;
         convergence threshold $\varepsilon > 0$}
\KwOutput{Optimal value function $V_\cI^*$;\; optimal policy $\pi_\cI^*$}
\KwParam{$\gamma \in [0,1)$}

\BlankLine
Initialise $V_0(s) \leftarrow 0$ for all $s \in \cS$\;
$k \leftarrow 0$\;

\Repeat{$\|V_{k+1} - V_k\|_\infty < \varepsilon$}{
  \ForEach{$s \in \cS$}{
    $\cA_\cI(s) \leftarrow \textsc{FeasibilityMonitor}(s,\, \cI(s),\, \delta_{\max})$\;
    \lIf{$\cA_\cI(s) = \textsc{Infeasible}$}{\textbf{halt} with infeasibility error}
    $V_{k+1}(s) \leftarrow \displaystyle\max_{a \in \cA_\cI(s)}
      \!\left[R(s,a) + \gamma \sum_{s'} \cP(s'\given s,a)\,V_k(s')\right]$\;
  }
  $k \leftarrow k + 1$\;
}

\ForEach{$s \in \cS$}{
  $\pi_\cI^*(s) \leftarrow \displaystyle\argmax_{a \in \cA_\cI(s)}
    \!\left[R(s,a) + \gamma \sum_{s'} \cP(s'\given s,a)\,V_\cI^*(s')\right]$\;
}

\textbf{return} $V_\cI^*,\; \pi_\cI^*$\;
\end{algorithm}

\subsection{Intent-Constrained Policy Iteration}

Algorithm~\ref{alg:icpi} constrains both the policy evaluation and
improvement steps to $\Pi_\cI$. No Lagrangian relaxation is required;
admissibility is guaranteed by construction.

\begin{algorithm}[H]
\caption{Intent-Constrained Policy Iteration (ICPI)}
\label{alg:icpi}
\KwInput{ICMDP $\cM_I$;\; convergence threshold $\varepsilon > 0$}
\KwOutput{Optimal admissible policy $\pi_\cI^*$}

\BlankLine
Initialise $\pi_0 \in \Pi_\cI$ arbitrarily
\Comment*[r]{any admissible initial policy}
$k \leftarrow 0$\;

\Repeat{$\pi_{k+1} = \pi_k$}{

  \Comment{--- Policy Evaluation ---}
  Solve $(I - \gamma P^{\pi_k}) V^{\pi_k} = R^{\pi_k}$ for $V^{\pi_k}$\;
  $Q_\cI^{\pi_k}(s,a) \leftarrow R(s,a) + \gamma \displaystyle\sum_{s'}
    \cP(s'\given s,a)\,V^{\pi_k}(s')
    \quad \forall s,\, a \in \cA_\cI(s)$\;

  \Comment{--- Policy Improvement (intent-filtered) ---}
  \ForEach{$s \in \cS$}{
    $\pi_{k+1}(s) \leftarrow \argmax_{a \in \cA_\cI(s)}\, Q_\cI^{\pi_k}(s,a)$\;
  }

  $k \leftarrow k + 1$\;
}

\textbf{return} $\pi_\cI^* \leftarrow \pi_k$\;
\end{algorithm}

% -----------------------------------------------------------------------
\section{Model-Free Reinforcement Learning}
\label{sec:model-free}
% -----------------------------------------------------------------------

\subsection{Masked Q-Learning}

Algorithm~\ref{alg:masked-q} estimates $Q_\cI^*$ from sampled transitions
without requiring explicit knowledge of $\cP$. The only modification to
standard Q-learning is that (i) updates are restricted to $a \in \cA_\cI(s)$
and (ii) the bootstrap target takes the max over $\cA_\cI(s')$.

\begin{algorithm}[H]
\caption{Masked Q-Learning}
\label{alg:masked-q}
\KwInput{Environment with state space $\cS$, action space $\cA$;\;
         intent-constraint function $\cI$;\;
         number of episodes $N$;\;
         episode horizon $T$}
\KwOutput{Estimated $Q_\cI^*$}
\KwParam{Learning rate $\alpha \in (0,1]$;\; discount $\gamma \in [0,1)$;\;
          exploration rate $\epsilon > 0$}

\BlankLine
Initialise $Q_\cI(s,a) \leftarrow 0$ for all $s \in \cS,\, a \in \cA$\;

\For{episode $= 1$ \KwTo $N$}{
  Observe initial state $s_0$\;
  \For{$t = 0$ \KwTo $T-1$}{
    $\cA_\cI(s_t) \leftarrow \textsc{FeasibilityMonitor}(s_t, \cI(s_t), \delta_{\max})$\;
    \lIf{$\cA_\cI(s_t) = \textsc{Infeasible}$}{\textbf{break}}

    \Comment{$\epsilon$-greedy over admissible actions only}
    \eIf{$\text{Uniform}(0,1) < \epsilon$}{
      $a_t \leftarrow \text{Uniform}(\cA_\cI(s_t))$\;
    }{
      $a_t \leftarrow \argmax_{a \in \cA_\cI(s_t)} Q_\cI(s_t, a)$\;
    }

    Execute $a_t$;\; observe $r_t,\, s_{t+1}$\;

    $\cA_\cI(s_{t+1}) \leftarrow \cI(s_{t+1})$\;

    \Comment{Masked Bellman update}
    $Q_\cI(s_t, a_t) \leftarrow Q_\cI(s_t, a_t)
      + \alpha \!\left[r_t + \gamma \max_{a' \in \cA_\cI(s_{t+1})}
        Q_\cI(s_{t+1}, a') - Q_\cI(s_t, a_t)\right]$\;

    $s_t \leftarrow s_{t+1}$\;
  }
}

\textbf{return} $Q_\cI$\;
\end{algorithm}

\subsection{Masked Policy Gradient}

Algorithm~\ref{alg:masked-pg} applies REINFORCE with a masked softmax
policy. Inadmissible actions receive logit $-\infty$ before the softmax,
guaranteeing zero probability mass on normatively excluded actions
\citep{huang2022} during both training and execution.

\begin{algorithm}[H]
\caption{Masked Policy Gradient (REINFORCE with Action Masking)}
\label{alg:masked-pg}
\KwInput{Policy network $\pi_\theta$ with parameters $\theta$;\;
         intent-constraint function $\cI$;\;
         number of episodes $N$;\;
         episode horizon $T$}
\KwOutput{Trained parameters $\theta^*$}
\KwParam{Learning rate $\eta > 0$;\; discount $\gamma \in [0,1)$}

\BlankLine
Initialise $\theta$ randomly\;

\For{episode $= 1$ \KwTo $N$}{
  Observe $s_0$;\; $\tau \leftarrow \{\}$\;

  \For{$t = 0$ \KwTo $T-1$}{
    $\cA_\cI(s_t) \leftarrow \textsc{FeasibilityMonitor}(s_t, \cI(s_t), \delta_{\max})$\;

    \Comment{Compute masked logits and softmax policy}
    $z_t(a) \leftarrow \begin{cases}
      \pi_\theta(a \given s_t) & \text{if } a \in \cA_\cI(s_t) \\
      -\infty                  & \text{otherwise}
    \end{cases}$\;

    $\hat{\pi}(a \given s_t) \leftarrow
      \dfrac{\exp(z_t(a))}{\sum_{a'} \exp(z_t(a'))}$\;

    Sample $a_t \sim \hat{\pi}(\cdot \given s_t)$;\;
    Execute $a_t$;\; observe $r_t,\, s_{t+1}$\;
    Append $(s_t, a_t, r_t)$ to $\tau$\;
    $s_t \leftarrow s_{t+1}$\;
  }

  \Comment{Compute discounted returns}
  \For{$t = 0$ \KwTo $T-1$}{
    $G_t \leftarrow \displaystyle\sum_{t'=t}^{T-1} \gamma^{t'-t} r_{t'}$\;
  }

  \Comment{Policy gradient update}
  $\theta \leftarrow \theta + \eta \displaystyle\sum_{t=0}^{T-1}
    G_t \nabla_\theta \log \hat{\pi}(a_t \given s_t)$\;
}

\textbf{return} $\theta^* \leftarrow \theta$\;
\end{algorithm}

% -----------------------------------------------------------------------
\section{Full Agentic Commerce Pipeline}
\label{sec:pipeline}
% -----------------------------------------------------------------------

Algorithm~\ref{alg:pipeline} integrates all components into the complete
end-to-end execution loop for an Agentic Commerce transaction. At each
time step the agent (i) updates its POMDP belief over hidden market state,
(ii) parses or re-evaluates intent constraints, (iii) checks feasibility,
(iv) selects the optimal admissible action, and (v) logs a verifiable
certificate of execution.

\begin{algorithm}[H]
\caption{Agentic Commerce ICMDP Execution Pipeline}
\label{alg:pipeline}
\KwInput{Natural language instruction $l$;\;
         pre-trained policy $\pi_\cI^*$ (or $Q_\cI^*$);\;
         initial belief $b_0 \in \Delta(\cS)$;\;
         transaction horizon $T$;\;
         observation model $\cZ$}
\KwOutput{Transaction log $\mathcal{T}$;\; execution certificate $\mathcal{C}$}

\BlankLine
$\mathcal{T} \leftarrow \{\}$;\;
$\mathcal{C} \leftarrow \{\}$\;

\Comment{Step 1 — Parse intent once at transaction start}
$\hat{s}_0 \leftarrow \argmax_{s} b_0(s)$\;
$\cA_\cI \leftarrow \textsc{LLMIntentParser}(l,\, \hat{s}_0,\, \cA,\, \theta)$\;
\Comment*[r]{Algorithm~\ref{alg:intent-parser}}

\For{$t = 0$ \KwTo $T - 1$}{

  \Comment{Step 2 — POMDP belief update}
  \lIf{$t > 0$}{$b_t \leftarrow \textsc{BeliefUpdate}(b_{t-1},\, a_{t-1},\, o_t,\, \cP,\, \cZ)$}

  \Comment{Step 3 — State estimate from belief}
  $\hat{s}_t \leftarrow \argmax_{s} b_t(s)$\;

  \Comment{Step 4 — Re-evaluate admissibility at current state}
  $\cA_\cI(s_t) \leftarrow \{a \in \cA_\cI : \text{all predicates satisfied at } \hat{s}_t\}$\;

  \Comment{Step 5 — Feasibility check}
  $\cA_\cI(s_t) \leftarrow \textsc{FeasibilityMonitor}(\hat{s}_t,\, \cA_\cI(s_t),\, \delta_{\max})$\;
  \Comment*[r]{Algorithm~\ref{alg:feasibility}}
  \lIf{$\cA_\cI(s_t) = \textsc{Infeasible}$}{\textbf{abort} transaction}

  \Comment{Step 6 — Select optimal admissible action}
  $a_t^* \leftarrow \pi_\cI^*(\hat{s}_t)$\Comment*[r]{or $\argmax_{a \in \cA_\cI(s_t)} Q_\cI^*(s_t, a)$}

  \Comment{Step 7 — Execute and observe}
  Execute $a_t^*$;\; receive reward $r_t$;\; observe $o_{t+1}$\;

  \Comment{Step 8 — Log certificate entry}
  $\mathcal{C} \leftarrow \mathcal{C} \cup
    \{(t,\; b_t,\; \cA_\cI(s_t),\; a_t^*,\; \text{proof: } a_t^* \in \cI(\hat{s}_t))\}$\;
  $\mathcal{T} \leftarrow \mathcal{T} \cup \{(t,\; a_t^*,\; r_t,\; o_{t+1})\}$\;
}

\textbf{return} $\mathcal{T},\; \mathcal{C}$\;
\end{algorithm}

\paragraph{Certificate of execution.}
Each entry in $\mathcal{C}$ records the belief state $b_t$, the
admissible set $\cA_\cI(s_t)$, the chosen action $a_t^*$, and a proof
that $a_t^* \in \cI(\hat{s}_t)$. This provides the user with a formal,
step-by-step audit trail that the agent operated within intent boundaries
throughout the transaction — a guarantee unavailable in reward-shaped or
expectation-constrained systems.

% -----------------------------------------------------------------------
\section{Complexity Notes}
\label{sec:complexity}
% -----------------------------------------------------------------------

\begin{itemize}[leftmargin=*]
  \item \textbf{ICVI} converges in $O\!\left(\frac{|\cS|^2 |\cA|}{(1-\gamma)\varepsilon}\right)$
  iterations \citep{puterman1994}, reduced to
  $O\!\left(\frac{|\cS|^2 |\cA_\cI|}{(1-\gamma)\varepsilon}\right)$ under
  intent masking when $|\cA_\cI(s)| \ll |\cA|$.

  \item \textbf{ICPI} converges in at most $|\cA_\cI|^{|\cS|}$ iterations
  (polynomial in practice).

  \item \textbf{Masked Q-Learning} inherits the $O(|\cS||\cA_\cI|)$ per-step
  cost of standard Q-learning, with reduced exploration space.

  \item \textbf{Intent parsing} (Algorithm~\ref{alg:intent-parser}) introduces
  a one-time LLM inference cost at transaction start, plus re-evaluation cost
  at each step where state-dependent predicates change.

  \item \textbf{Feasibility monitoring} is $O(|\cA|)$ per state and should be
  invoked before every policy query.
\end{itemize}

% -----------------------------------------------------------------------
%  Bibliography
% -----------------------------------------------------------------------
\bibliographystyle{plainnat}
\bibliography{references}

\end{document}
